\documentclass[12pt,a4paper]{article}

\usepackage[margin=2.5cm]{geometry}
\usepackage{booktabs}
\usepackage{array}
\usepackage{hyperref}
\usepackage{parskip}

\title{The Degu (\textit{Octodon degus}): A Comprehensive Overview}
\author{A. Researcher}
\date{\today}

\begin{document}

\maketitle
\tableofcontents
\newpage

\section{Introduction}

The degu (\textit{Octodon degus}) is a small caviomorph rodent native to the western slopes of the Andes mountains in Chile.
Often described as a cross between a rat and a squirrel, the degu has attracted considerable scientific interest, e.g., in the fields of diabetes research, neuroscience, and chronobiology.
The animal belongs to the family Octodontidae, which includes several other burrowing rodents found in South America, e.g., the rock rat (\textit{Aconaemys fuscus}) and the viscacha rat (\textit{Octomys mimax}).

In recent decades the degu has become an increasingly popular laboratory animal, i.e., it is now routinely used as a model organism for studying type 2 diabetes, Alzheimer's disease, and circadian rhythms.
This popularity stems from a number of physiological traits that make it particularly useful, e.g., its inability to handle dietary sucrose without developing diabetic symptoms, and its naturally diurnal activity pattern, which is unusual among rodents.
This document provides an overview of the biology, ecology, behaviour, and husbandry of the degu, etc., drawing on current literature.

\section{Taxonomy and Phylogeny}

The degu was first described by the naturalist Rodolfo Amando Philippi in 1860, though it had been known to indigenous Chilean peoples long before formal taxonomic description.
The genus name \textit{Octodon} refers to the figure-eight shape formed by the occlusal surfaces of the animal's molar teeth, i.e., each molar has two circular enamel loops joined together, resembling the numeral eight.
This dental morphology is a diagnostic trait for the family Octodontidae.

The family Octodontidae is part of the superfamily Octodontoidea within the suborder Hystricomorpha.
Caviomorph rodents are found exclusively in the New World, and are thought to have arrived in South America approximately 37--40 million years ago, possibly by rafting across the Atlantic from Africa.
The phylogenetic relationships within the Octodontidae are still debated, e.g., some authorities place the tuco-tuco (\textit{Ctenomys} spp.)
within the Octodontidae while others assign it to a separate family, i.e., Ctenomyidae.

Table~\ref{tab:taxonomy} summarises the full taxonomic classification of the degu.

\begin{table}[ht]
\centering
\caption{Taxonomic classification of \textit{Octodon degus}.}
\label{tab:taxonomy}
\begin{tabular}{ll}
\toprule
\textbf{Rank} & \textbf{Name} \\
\midrule
Kingdom  & Animalia \\
Phylum   & Chordata \\
Class    & Mammalia \\
Order    & Rodentia \\
Suborder & Hystricomorpha \\
Family   & Octodontidae \\
Genus    & \textit{Octodon} \\
Species  & \textit{O.\ degus} \\
\bottomrule
\end{tabular}
\end{table}

\section{Morphology}

Adult degus typically weigh between 170 and 300 g, with a body length of approximately 25--31 cm including the tail.
The tail, which accounts for roughly one third of total body length, is notably different from that of most rodents in that it can be shed (i.e., autotomised) when grabbed by a predator.
Unlike some lizards, however, the degu cannot regenerate a lost tail, so this is a one-time escape mechanism.

The fur of the degu is generally ochraceous-brown on the dorsal surface, fading to a pale yellowish-cream on the ventral surface.
The tail terminates in a characteristic tuft of dark hair, which is visible even in young animals.
The eyes are large and laterally placed, providing a wide field of view, which is an important adaptation for a prey species living in open habitats.
The ears are moderately large and are used both for hearing and thermoregulation, e.g., excess body heat can be dissipated through the large network of blood vessels in the pinnae.

The dental formula of the degu is: incisors 1/1, canines 0/0, premolars 1/1, molars 3/3, giving a total of 20 teeth.
The incisors are unpigmented (white), which distinguishes the degu from many other hystricomorph rodents, e.g., the guinea pig (\textit{Cavia porcellus}), which has orange-pigmented incisors.
The figure-eight shaped molars, as mentioned in the previous section, are hypsodont (i.e., high-crowned) and are continuously worn down by the animal's abrasive diet of grasses and other fibrous plant material.

\subsection{Sexual Dimorphism}

Degus show moderate sexual dimorphism.
Males are generally slightly larger than females, i.e., they tend to be heavier by approximately 10--15\% on average.
There are no dramatic external differences between the sexes aside from the obvious anatomical differences, e.g., females have four pairs of mammary glands.
Sexing of young animals can be challenging, and is typically done by measuring the anogenital distance, which is greater in males.

\section{Distribution and Habitat}

\textit{Octodon degus} is endemic to Chile, where it occupies a fairly narrow range along the western slopes of the Andes, primarily in the regions of Atacama, Coquimbo, Valparaiso, Metropolitana, and O'Higgins.
The species is found predominantly at elevations between sea level and approximately 1,200 m, though some populations have been recorded at higher altitudes, e.g., up to 3,000 m in the Andes.

The preferred habitat of the degu is semi-arid scrubland (i.e., the Chilean \textit{matorral}), characterised by a sparse cover of thorny shrubs, cacti, and drought-resistant grasses.
The soil in these habitats is typically rocky or sandy, which facilitates the construction of the burrow systems that degus rely upon for shelter, predator avoidance, and thermoregulation.
Degus tend to avoid both dense forest and open desert, preferring the transitional habitat in between, e.g., hillsides with patchy vegetation cover.

\section{Behaviour}

\subsection{Social Organisation}

Degus are highly social animals that live in family groups typically consisting of 2--9 individuals, though larger aggregations have been recorded under favourable conditions.
The social group shares a communal burrow system, which may be quite extensive, e.g., comprising dozens of chambers and tunnels extending several metres below the surface.
Burrow construction is a cooperative activity, with all members of the group contributing to digging, maintenance, and debris clearance.

Within the group, degus establish a dominance hierarchy, i.e., a linear rank order in which higher-ranking individuals have priority access to food, mates, and preferred resting sites.
Dominance is maintained through a range of behaviours including scent marking, vocalisation, and physical aggression, e.g., boxing, chasing, and fur-pulling.
Despite this hierarchy, outright aggression is relatively rare under natural conditions, and the group is cohesive, with individuals frequently engaging in allogrooming.

\subsection{Communication}

Degus are remarkably vocal for rodents and possess a large repertoire of calls, i.e., researchers have identified at least 15 distinct vocalisations in wild populations, ranging from contact calls used to maintain group cohesion to alarm calls that warn of aerial predators (e.g., hawks and owls).
Degus also communicate via ultrasonic vocalisations that are inaudible to humans, which are used particularly during mother--pup interactions.

Olfactory communication is also important, e.g., degus possess well-developed scent glands on the flank, belly, and chin, which are used to mark territory, mates, and burrow entrances.
Urine washing, i.e., the behaviour of rubbing urine onto the palmar surface of the hind feet, is commonly observed and may serve a communicative function.

\subsection{Activity Patterns}

One of the most distinctive and scientifically valuable traits of the degu is its diurnal activity pattern.
Most rodents used in laboratory research, e.g., mice and rats, are nocturnal, which creates complications when extrapolating findings to humans.
Degus, by contrast, are active primarily during the day, i.e., their circadian rhythms are more similar to those of humans than those of the standard rodent models.
This makes them particularly valuable for chronobiological research, e.g., the study of circadian disruption in shift workers or patients with sleep disorders.

\section{Diet and Foraging}

Degus are primarily herbivorous, subsisting on a diet of grasses, leaves, seeds, and bark.
They are particularly fond of the leaves and bark of Chilean shrubs such as \textit{Puya} spp.
and \textit{Opuntia} cacti.
During periods of food scarcity, e.g., the dry season, they will also consume roots and bulbs excavated from the soil.
Like many arid-land rodents, degus are efficient in their use of dietary water, i.e., they can obtain much of their water requirement from their food and produce concentrated urine to reduce water loss.

Foraging typically occurs in short bursts near the burrow entrance, with individuals rarely venturing more than 30--50 m from the nearest refuge.
This conservative foraging strategy reduces exposure to predators, e.g., raptors, foxes, and culpeos (\textit{Lycalopex culpaeus}), at the cost of limiting access to more distant food patches.
Degus will sometimes cache food near their burrows, particularly seeds and dry plant material, which they may consume during periods of inclement weather.

Table~\ref{tab:diet} shows representative nutritional components of foods commonly consumed by degus in the wild, and Table~\ref{tab:lab_diet} compares recommended captive diets.

\begin{table}[ht]
\centering
\caption{Approximate nutritional composition of selected degu food items (per 100 g dry matter).}
\label{tab:diet}
\begin{tabular}{lrrrr}
\toprule
\textbf{Food item} & \textbf{Protein (\%)} & \textbf{Fat (\%)} & \textbf{Fibre (\%)} & \textbf{Sucrose (\%)} \\
\midrule
Timothy grass hay    & 8.5  & 2.1  & 32.0 & 0.4 \\
\textit{Opuntia} pad & 4.2  & 0.6  & 18.5 & 1.1 \\
Grass seeds (mixed)  & 12.8 & 5.7  & 10.2 & 2.3 \\
Bark (\textit{Puya}) & 3.1  & 0.8  & 45.3 & 0.3 \\
\bottomrule
\end{tabular}
\end{table}

\begin{table}[ht]
\centering
\caption{Comparison of captive degu diet recommendations from selected authorities.}
\label{tab:lab_diet}
\begin{tabular}{p{4cm}p{4.5cm}p{4.5cm}}
\toprule
\textbf{Authority} & \textbf{Recommended diet} & \textbf{Foods to avoid} \\
\midrule
Correa et al.\ (2006) & High-fibre hay ad libitum; rodent pellets sparingly & All sugars, i.e., fruit, honey, molasses \\
Jekl \& Redrobe (2011) & Chinchilla pellets, grass hay, limited fresh greens & Sucrose, fructose, e.g., apple, banana \\
ACLAM (2020) & Specialist degu pellets; unlimited timothy hay & High-fat seeds, e.g., sunflower; any sweeteners \\
\bottomrule
\end{tabular}
\end{table}

\section{Reproduction}

Degus are seasonal breeders in the wild, with the main breeding season occurring in the austral winter and spring (i.e., June--November), timed such that pups are born when vegetation growth is most abundant.
However, captive animals may breed year-round if maintained under stable photoperiod and temperature conditions.

The gestation period of the degu is approximately 90 days, which is long relative to body size and reflects the precocial development of the pups.
Litter size typically ranges from 1 to 12, with an average of approximately 5--6 pups.
Neonates are born with a full coat of fur, open eyes, and the ability to move around the nest almost immediately, i.e., they are highly precocial compared to altricial rodents such as mice and rats.
Birth weight is approximately 13--16 g.

Pups are nursed by the mother for approximately 4--6 weeks, though they begin consuming solid food within the first week of life.
Weaning is a gradual process, e.g., the mother does not abruptly cease nursing but rather gradually reduces the frequency of nursing bouts as the pups become more independent.
Allomaternal care is common in group-living degus, i.e., females other than the biological mother will nurse, groom, and protect the pups.

Sexual maturity is reached at approximately 12--16 weeks of age in females and slightly later in males.
Life expectancy in captivity is typically 6--8 years, though individuals surviving to 10 years have been recorded.

\section{Medical Significance and Use as a Model Organism}

\subsection{Diabetes Research}

The degu's susceptibility to diet-induced diabetes mellitus is one of its most important characteristics from a biomedical perspective, i.e., it is one of the few rodent species in which type 2 diabetes can be reliably induced by dietary manipulation alone.
Specifically, feeding degus diets containing sucrose or fructose leads to hyperglycaemia, cataracts, and pancreatic changes closely resembling those seen in human type 2 diabetes.
This occurs because degus lack sufficient hepatic fructokinase activity to metabolise dietary fructose, e.g., consumption of as little as 5\% sucrose in the diet can produce overt diabetes within a few months.

\subsection{Neuroscience and Alzheimer's Disease}

Degus have attracted interest as a model for Alzheimer's disease (AD) research.
Aged degus spontaneously develop amyloid plaques and neurofibrillary tangles in the cerebral cortex and hippocampus, i.e., the two hallmark pathological features of AD, without genetic modification.
This is in contrast to most other rodent models of AD, e.g., transgenic mice, which require artificial overexpression of human amyloid precursor protein or presenilin genes to develop AD-like pathology.
The cognitive deficits associated with ageing in degus, e.g., impaired spatial learning and memory, are also broadly similar to those seen in early-stage AD patients.

\subsection{Chronobiology}

As noted above, the diurnal activity pattern of the degu makes it valuable for studying circadian rhythms, e.g., the effects of light pollution, shift work, and sleep deprivation on human health.
The degu's circadian clock has been well characterised at both the behavioural and molecular level, i.e., the master pacemaker in the suprachiasmatic nucleus (SCN) of the hypothalamus has been shown to function similarly to that of humans in terms of its response to light and social cues.

\section{Captive Husbandry}

Degus have become popular as pets in many countries, i.e., they are kept throughout Europe, North America, and Japan.
They are intelligent, curious, and active, which makes them engaging to observe, but they also have specific husbandry requirements that must be met to ensure their health and welfare, e.g., they require large enclosures, dust baths, and companionship.

\subsection{Housing}

Degus should never be housed alone, i.e., solitary housing causes severe psychological distress in these highly social animals.
Pairs or small groups are recommended for pet degus.
The enclosure should be as large as practical, with a minimum floor area of approximately 100 cm $\times$ 50 cm for a pair, and should include multiple levels since degus are agile climbers.
Wire-mesh cages are generally preferred over glass aquaria, e.g., because they provide better ventilation and allow the animals to climb the walls.
Solid-floored enclosures are preferable to wire-bottomed cages, which can cause pododermatitis (i.e., inflammation of the feet).

Table~\ref{tab:housing} provides a summary of recommended housing parameters for captive degus.

\begin{table}[ht]
\centering
\caption{Recommended husbandry parameters for captive degus.}
\label{tab:housing}
\begin{tabular}{ll}
\toprule
\textbf{Parameter} & \textbf{Recommended value} \\
\midrule
Minimum enclosure floor area (pair) & 100 cm $\times$ 50 cm \\
Minimum enclosure height           & 60 cm \\
Ambient temperature                & 18--22\,\textdegree C \\
Relative humidity                  & 40--60\,\% \\
Photoperiod                        & 12L:12D (i.e., 12 h light, 12 h dark) \\
Bedding depth                      & $\geq$5 cm (e.g., aspen shavings, paper bedding) \\
Dust bath frequency                & 2--3 times per week \\
Group size                         & $\geq$2 (never house alone) \\
\bottomrule
\end{tabular}
\end{table}

\subsection{Diet in Captivity}

As discussed in Section 6, the degu's inability to tolerate dietary sugars requires careful attention to feeding.
A high-fibre diet based on timothy grass hay, with small amounts of specialist degu pellets, is the generally accepted standard.
Fresh vegetables may be offered in small quantities, e.g., leafy greens such as kale or romaine lettuce, but fruit should be strictly avoided due to its sugar content.
Seeds and nuts are also high in fat and should be offered sparingly, e.g., as occasional enrichment items rather than dietary staples.

\subsection{Enrichment}

Degus are highly intelligent and require substantial environmental enrichment to prevent boredom and stereotypic behaviours, e.g., bar-chewing and repetitive pacing.
Enrichment items may include cardboard tubes, untreated wooden blocks for gnawing, climbing branches, ropes, and exercise wheels.
An exercise wheel, i.e., a solid-surfaced wheel with a diameter of at least 25 cm, is considered essential, as degus in the wild may travel considerable distances each day.
Dust baths (using chinchilla sand, not clay-based dust) should be offered several times per week, as they are used for coat maintenance and appear to have a positive psychological effect.

\section{Conservation Status}

\textit{Octodon degus} is currently listed as Least Concern on the IUCN Red List, i.e., it is not considered to be at risk of extinction in the near future.
The species is abundant within its native range and has proven highly adaptable, e.g., it has colonised disturbed habitats such as roadsides and agricultural margins.
Some localised declines have been reported in association with urban expansion around Santiago, but the overall population trend is considered stable.

Paradoxically, the degu has become an invasive species in some regions outside its native range, e.g., populations of escaped or released pet degus have established themselves in parts of Germany and the United Kingdom.
These feral populations are considered potentially problematic, i.e., they could compete with native rodents, e.g., the bank vole (\textit{Myodes glareolus}) and wood mouse (\textit{Apodemus sylvaticus}), though the extent of any such competition has not been rigorously quantified.

\section{Conclusion}

The degu (\textit{Octodon degus}) is a remarkable rodent that has earned a distinguished place in both biomedical research and the exotic pet trade.
Its unique combination of traits, e.g., diurnal activity, high sociality, susceptibility to diabetes, and spontaneous Alzheimer's-like pathology, makes it one of the most scientifically valuable small mammals available to researchers.
At the same time, its intelligence, engaging behaviour, and relatively long lifespan make it an appealing companion animal for those willing to meet its husbandry requirements.

As with all exotic pets, prospective owners should research thoroughly before acquiring degus, i.e., they should be aware of the species' social needs, dietary restrictions, and the financial commitment required for veterinary care.
Specialist exotic animal veterinarians should be consulted regularly, e.g., for annual health checks and dental examinations, since dental disease is among the most common health problems in captive degus.

Future research directions include further elucidation of the degu's circadian biology, development of improved AD models using aged degus, and conservation genetics of wild populations, etc.
Given the breadth of scientific and popular interest in this small Chilean rodent, it seems likely that the degu will continue to attract attention for many years to come.

\end{document}
